\chapter{Surface and Interface Physics}

At a surface, the three-dimensional periodicity of a crystal is broken. This broken symmetry gives rise to surface-specific electronic states, reconstruction, and entirely new physics in reduced dimensions---most spectacularly, the quantum Hall effect.

\section{Surface Crystallography}

A clean crystal surface reconstructs to minimize its energy. The Si(111)-$7\times7$ reconstruction (Chapter~1) is the most complex known surface structure, with 49 atoms in the surface unit cell.

\section{The Quantum Hall Effect}

\subsection{Integer Quantum Hall Effect (1980)}

Klaus von Klitzing~\cite{klitzing1986qhe} discovered that the Hall resistance of a two-dimensional electron gas (2DEG) in a GaAs/AlGaAs heterostructure, at low temperatures and high magnetic fields, is quantized:
\begin{equation}
  R_H = \frac{h}{ne^2} = \frac{\SI{25812.807}{\ohm}}{n}, \quad n = 1, 2, 3, \ldots
\end{equation}

The quantization is exact to better than 1 part in $10^{9}$, making it the basis of the international resistance standard. Von Klitzing received the 1985 Nobel Prize.

\subsection{Fractional Quantum Hall Effect (1982)}

Tsui, Stormer, and Gossard discovered plateaus at fractional filling factors $\nu = 1/3, 2/5, 3/7, \ldots$ in ultra-high-mobility GaAs samples. Laughlin explained these as an incompressible quantum liquid of electrons, for which all three received the 1998 Nobel Prize.

The $\nu = 1/3$ state has quasiparticles with fractional charge $e/3$---experimentally confirmed by shot noise measurements.


\section{p--n Junctions}

The p--n junction is the building block of semiconductor electronics. The current--voltage characteristic is:
\begin{equation}
  I = I_0\left(e^{eV/k_BT} - 1\right)
\end{equation}

\begin{table}[H]
\centering
\caption{Parameters of selected p--n junction devices.}
\label{tab:pn}
\begin{tabular}{@{}llSl@{}}
\toprule
Device & Material & {Turn-on voltage (\si{\volt})} & Application \\
\midrule
Rectifier diode & Si     & 0.7  & Power conversion \\
Schottky diode  & GaAs   & 0.3  & High-speed switching \\
LED (red)       & GaAsP  & 1.8  & Indicator lights \\
LED (blue)      & InGaN  & 3.0  & Displays, lighting \\
Solar cell      & Si     & 0.5  & Photovoltaics \\
Laser diode     & GaAs   & 1.4  & Telecommunications \\
\bottomrule
\end{tabular}
\end{table}

\section{Light-Emitting Diodes}

When a forward-biased p--n junction recombines electrons and holes, photons are emitted with energy $\approx E_g$. The LED spectrum peaks at $\lambda \approx hc/E_g$:

\begin{table}[H]
\centering
\caption{LED emission wavelengths and corresponding band gaps.}
\label{tab:led}
\begin{tabular}{@{}llSS@{}}
\toprule
Color & Material & {$\lambda$ (\si{\nano\metre})} & {$E_g$ (\si{\electronvolt})} \\
\midrule
Infrared & GaAs      & 940  & 1.32 \\
Red      & GaAsP     & 660  & 1.88 \\
Orange   & GaAsP     & 610  & 2.03 \\
Green    & GaP:N     & 565  & 2.20 \\
Blue     & InGaN     & 470  & 2.64 \\
UV       & AlGaN     & 340  & 3.65 \\
\bottomrule
\end{tabular}
\end{table}

The development of efficient blue LEDs by Akasaki, Amano, and Nakamura (Nobel Prize 2014) enabled white LED lighting, which is now the dominant lighting technology worldwide.

\vspace{1cm}
\noindent\rule{\textwidth}{0.4pt}
\textit{The physics of structures with dimensions comparable to the electron wavelength---nanostructures---is the subject of the next chapter.}
